\documentclass[12pt]{article}
\usepackage[utf8]{inputenc}
\usepackage[T1]{fontenc}
\usepackage{hyperref}
\title{Wotjobaluk Language Resources for K-12 Curriculum Development}
\author{Compiled by Claude Sonnet 5.5 (claude-sonnet-5-5) for Andr\'as Kornai}
\date{}
\begin{document}
\maketitle

\section{Introduction}
Wotjobaluk (ISO 639-3 xwt) is an Australian language of the Wimmera region of Victoria. Sources treat it as a dialect of Wergaia, which has its own ISO code (weg), so most resources below are filed under Wergaia. The language has no native speakers, and a community revival effort exists. The goal is a heritage and revitalization curriculum.

\section{Primary Educational Resources}
\begin{itemize}
\item A 20-week Wergaia language program was launched in mid-2021 at the Wotjobaluk Knowledge Place in Dimboola, as described on the Wikipedia page for the Wotjobaluk people. Its current status and materials were not confirmed: \url{https://en.wikipedia.org/wiki/Wotjobaluk_people}
\end{itemize}

\section{Linguistic and Reference Materials}
\begin{itemize}
\item Reid, Julie (2007), \emph{Wergaia Community Grammar and Dictionary}. It is a consultation draft with a sketch grammar and a revised wordlist, and is not intended as a language course. The search result was not the dictionary itself, so check the Victorian Collections listing: \url{https://victoriancollections.net.au/collections?q=wergaia}
\item Chapter ``The rebirth of Wergaia: a collaborative effort'', Sydney University Press: \url{https://open.sydneyuniversitypress.com.au/9781920899554/rl-the-rebirth-of-wergaia-a-collaborative-effort.html}
\item Wikipedia, Wergaia language: \url{https://en.wikipedia.org/wiki/Wergaia_language}
\item Wikipedia, Wergaia dialect: \url{https://en.wikipedia.org/wiki/Wergaia_dialect}
\item Endangered Languages Project entry: \url{https://endangeredlanguages.com/node/97779}
\end{itemize}

\section{Religious and Community Resources}
\begin{itemize}
\item South Australian Museum archive page on the Wotjobaluk language group: \url{https://www.samuseum.sa.gov.au/collection/archives/language_groups/wotjobaluk}
\end{itemize}

\section{Digital and Audio Resources}
No audio, video or software resources were verified in the searches.

\section{Curriculum Development Context}
There is a community revival effort and a draft grammar and dictionary, which is more than exists for many extinct languages. The draft is not a course, so lesson materials would need to be built with the Wotjobaluk community, drawing on the Wergaia work.
\end{document}